\chapter{Clase Taller 6}

\section{Notación escalar, vectorial y tensorial}

Un escalar se representa mediante un símbolo como $\phi$. Un vector se expresa en términos de sus componentes y de los vectores de la base:
\[
\vec A=A_i\,\hat e_i,
\qquad
\{A_i\}=\begin{pmatrix}A_1\\A_2\\A_3\end{pmatrix}.
\]
Para un tensor de segundo orden, la representación correspondiente es
\[
\mathbb{T}=T_{ij}\,\hat e_i\otimes\hat e_j,
\]
donde $\otimes$ denota el producto tensorial y $T_{ij}$ son las componentes del tensor.

\subsection{Convención de suma de Einstein}

Un índice que aparece dos veces en un término se suma; un índice que aparece una sola vez es libre y debe coincidir en ambos miembros de la igualdad. Por ejemplo,
\[
A_iB_i=\sum_{i=1}^{3}A_iB_i,
\qquad
C_i=T_{ij}A_j.
\]
En la segunda expresión, $i$ es libre y $j$ es un índice de suma. Los índices de suma pueden renombrarse sin alterar el resultado, siempre que se eviten conflictos con otros índices del término.

\section{Traza y contracciones}

\subsection{Traza de un tensor}

La traza es la suma de las componentes diagonales:
\[
\mathrm{Tr}\,\mathbb{S}=S_{11}+S_{22}+S_{33}=S_{ii}.
\]

\subsection{Contracciones dobles}

Es importante distinguir las siguientes contracciones:
\[
\begin{aligned}
S_{ij}S_{ij}
&=S_{11}^2+S_{12}^2+S_{13}^2\\
&\quad+S_{21}^2+S_{22}^2+S_{23}^2\\
&\quad+S_{31}^2+S_{32}^2+S_{33}^2
=\mathrm{Tr}(\mathbb{S}^T\mathbb{S}),
\end{aligned}
\]
\[
\begin{aligned}
S_{ij}S_{ji}
&=S_{11}^2+S_{12}S_{21}+S_{13}S_{31}\\
&\quad+S_{21}S_{12}+S_{22}^2+S_{23}S_{32}\\
&\quad+S_{31}S_{13}+S_{32}S_{23}+S_{33}^2
=\mathrm{Tr}(\mathbb{S}^2).
\end{aligned}
\]
El cambio de nombre de los índices da $S_{jk}S_{kj}=S_{ij}S_{ji}$. Las dos contracciones anteriores coinciden cuando $\mathbb{S}$ es simétrico.

\subsection{Forma cuadrática}

La contracción de un tensor con dos copias de un vector produce un escalar:
\[
\begin{aligned}
\vec a\cdot\mathbb{S}\cdot\vec a=S_{mn}a_ma_n
&=S_{11}a_1^2+S_{12}a_1a_2+S_{13}a_1a_3\\
&\quad+S_{21}a_2a_1+S_{22}a_2^2+S_{23}a_2a_3\\
&\quad+S_{31}a_3a_1+S_{32}a_3a_2+S_{33}a_3^2.
\end{aligned}
\]

\section{Tensores simétricos y antisimétricos}

Un tensor $\mathbb{S}$ es simétrico si $S_{ij}=S_{ji}$, mientras que un tensor $\mathbb{A}$ es antisimétrico si $A_{ij}=-A_{ji}$. En este último caso cada componente diagonal es nula:
\[
A_{11}=A_{22}=A_{33}=0.
\]
La contracción entre un tensor simétrico y uno antisimétrico se anula. Renombrando los índices de suma,
\[
S_{ij}A_{ij}=S_{ji}A_{ji}=-S_{ij}A_{ij}
\quad\Longrightarrow\quad S_{ij}A_{ij}=0.
\]

\section{Producto vectorial y determinante}

El símbolo de Levi-Civita $\varepsilon_{ijk}$ vale $+1$ para las permutaciones pares de $(1,2,3)$, $-1$ para las impares y cero si algún índice se repite.

\subsection{Producto vectorial y producto mixto}

En una base cartesiana de orientación positiva,
\[
(\vec A\times\vec B)_i=\varepsilon_{ijk}A_jB_k,
\qquad
(\vec A\times\vec B)\cdot\vec C=\varepsilon_{ijk}A_iB_jC_k.
\]
La magnitud del producto vectorial satisface
\[
(\vec A\times\vec B)\cdot(\vec A\times\vec B)
=|\vec A|^2|\vec B|^2-(\vec A\cdot\vec B)^2.
\]

\subsection{Determinante de una matriz}

Para una matriz $\mathbb{A}$ de dimensión $3\times3$,
\[
\begin{aligned}
\det\mathbb{A}&=\varepsilon_{ijk}A_{1i}A_{2j}A_{3k}\\
&=A_{11}(A_{22}A_{33}-A_{23}A_{32})\\
&\quad-A_{12}(A_{21}A_{33}-A_{23}A_{31})\\
&\quad+A_{13}(A_{21}A_{32}-A_{22}A_{31}).
\end{aligned}
\]

\section{Derivación en notación indicial}

La identidad $\partial x_i/\partial x_k=\delta_{ik}$ permite aplicar la regla del producto:
\[
\frac{\partial(x_ix_j)}{\partial x_k}
=\delta_{ik}x_j+x_i\delta_{jk}.
\]
Para un campo tensorial,
\[
\frac{\partial(T_{ij}x_j)}{\partial x_k}
=\frac{\partial T_{ij}}{\partial x_k}x_j+T_{ik}.
\]
Si $T_{ij}$ es constante, el primer término se anula y se recupera $\partial(T_{ij}x_j)/\partial x_k=T_{ik}$.

Las segundas derivadas de un campo escalar definen una forma cuadrática en los incrementos cartesianos:
\[
d^2\phi=\frac{\partial^2\phi}{\partial x_i\,\partial x_j}\,dx_i\,dx_j.
\]
Cada índice de suma aparece exactamente dos veces en este término.

\section{Coordenadas e invariantes}

\subsection{Transformaciones entre bases cartesianas}

Sea $\mathbb{Q}=\{q_{ij}\}$ una matriz ortogonal de cambio de base. Sus componentes satisfacen
\[
q_{ik}q_{jk}=\delta_{ij},
\qquad q_{ki}q_{kj}=\delta_{ij},
\qquad \det\mathbb{Q}=\pm1.
\]
Las componentes de un tensor se transforman según
\[
T'_{ij}=q_{ik}q_{jl}T_{kl},
\qquad \mathbb{T}'=\mathbb{Q}\mathbb{T}\mathbb{Q}^T.
\]
La traza y el determinante son invariantes bajo este cambio de base:
\[
\begin{aligned}
\mathrm{Tr}\,\mathbb{T}'
&=q_{ik}q_{il}T_{kl}
=\delta_{kl}T_{kl}=T_{kk}=\mathrm{Tr}\,\mathbb{T},\\
\det\mathbb{T}'
&=\det\mathbb{Q}\,\det\mathbb{T}\,\det\mathbb{Q}^T
=(\det\mathbb{Q})^2\det\mathbb{T}=\det\mathbb{T}.
\end{aligned}
\]

\subsection{Elemento de línea y coordenadas curvilíneas}

En coordenadas cartesianas, la base es fija:
\[
\vec r=x_i\hat e_i,
\qquad d\vec r=dx_i\hat e_i,
\qquad ds^2=dx_i\,dx_i.
\]
Si se expresa un vector en una base variable como $\vec r=r_i\hat e_i$, su diferencial incluye la variación de la base:
\[
d\vec r=dr_i\,\hat e_i+r_i\,d\hat e_i.
\]
En coordenadas curvilíneas $\xi_i$, se escribe $\vec r=\vec r(\xi_1,\xi_2,\xi_3)$ y se define la base coordenada $\vec g_i=\partial\vec r/\partial\xi_i$. Entonces,
\[
d\vec r=\vec g_i\,d\xi_i,
\qquad ds^2=g_{ij}\,d\xi_i\,d\xi_j,
\qquad g_{ij}=\vec g_i\cdot\vec g_j.
\]
En el caso cartesiano ortonormal, $g_{ij}=\delta_{ij}$ y se recupera la expresión anterior del elemento de línea.
